% ---------------------------------------------------------------------------
% Design decisions and scenario analysis: the funding-event model argued end to
% end. Synthetic experiments (exp19--exp28) are cited as calibrated evidence for
% the scenario map, not as performance claims; derivations are in MATH_NOTES.md
% (cited M1--M8). Input from complete_account.tex.
% ---------------------------------------------------------------------------
\section{Design decisions: modeling funding events, argued end to end}
\label{sec:design}

We have deliberately not touched real data yet. This section is therefore
written the only honest way a pre-data chapter can be: as a \emph{modeling
argument} --- what a funding event is statistically, why the architecture
follows from the data-generating reality rather than from taste, and, most
importantly, a \emph{scenario map}: in which worlds we expect this machinery to
work, in which worlds specific parameters stop being identifiable, and what we
predict will happen on real data before we see it. The synthetic campaign
(\code{exp19}--\code{exp28}) plays the role of calibrated intuition: each
scenario claim below was rehearsed in a world where the truth is known, so the
expectations come with a mechanism, not just a hunch. Derivations are collected
in \code{MATH\_NOTES.md} (M1--M8).

% ---------------------------------------------------------------------
\subsection{The data reality that dictates the architecture}
\label{sec:design_data}

Four properties of startup-funding data are terrain, not choices.

\begin{enumerate}[leftmargin=1.6em]
\item \textbf{Event-time-only updates.} A firm's covariates refresh only at its
own deals; between deals nothing is observed. Features are
last-observation-carried-forward by construction. Any model that pretends to
watch firms continuously is fiction; any feature pipeline that reads a value
recorded after the event is leakage.
\item \textbf{Positive-only labels.} We observe fundings, never rejections or
failed raises. There is no firm-level negative class, so the natural likelihood
is conditional on the event (``given a deal in sector $s$ at $t$, which
firm?''), not a pointwise classifier over firm-weeks (M8).
\item \textbf{A churning, partially observed universe.} Firms are founded,
raise for the first time (before which they are invisible), and leave by
failure or graduation (IPO/M\&A) --- mostly without labels. The candidate set
is dynamic and incomplete; a fraction of deals belongs to firms outside any
tracked universe.
\item \textbf{Resolution is a modeling choice.} Deal dates exist at day
resolution; weekly bucketing is something \emph{we} impose. What aggregation
destroys is quantifiable and is part of the scenario map
(\S\ref{sec:design_scenarios}).
\end{enumerate}

% ---------------------------------------------------------------------
\subsection{The two stages are a factorization, not an approximation}
\label{sec:design_factorization}

The cornerstone is the marked-point-process factorization (Daley--Vere-Jones;
M1): the firm-level intensity always decomposes as
\begin{equation}
\lambda_i(t)\;=\;\underbrace{\Lambda_{s(i)}(t)}_{\text{ground: when/where}}
\;\cdot\;
\underbrace{\Prob\!\left(i \mid s(i),\,t,\,\cH_t\right)}_{\text{mark: which firm}}.
\label{eq:design_factorization}
\end{equation}
This is an identity, so ``two-stage'' loses nothing by construction. All
modeling content lives in how each factor is parametrized, and those
restrictions are named exhaustively in M1 ([G1]--[G6] for the ground factor,
[P1]--[P6] for the mark factor): log-linear weekly autoregression or
exponential-kernel event-time Hawkes for $\Lambda_s$; softmax over the live
risk set with linear utilities, one cooldown covariate, and an outside option
for the mark. Two structural consequences follow (M8): the joint likelihood
separates whenever the factors share no parameters, so the mark model is a
valid partial likelihood \emph{without any timing model at all}; and the places
where we do couple history (cooldown features, the funded-pool rule) remain
valid because they condition on $\cH_t$ only.

The value of stating it this way is rhetorical discipline: whenever a reviewer
asks ``why two stages?'', the answer is that the decomposition is exact and the
\emph{restrictions inside each factor} are the discussable content --- each of
which has either a derivation or a standing falsification test
(\S\ref{sec:design_bench}).

% ---------------------------------------------------------------------
\subsection{Parametrizing the ground factor: sector Hawkes, held loosely}
\label{sec:design_ground}

\paragraph{Sign structure decides the level.} Sector-level contagion is
positive (a deal raises attention and liquidity for neighbors); the firm-level
self-effect is negative (a firm that just raised leaves the market for months).
Nonnegative Hawkes kernels can express the former and \emph{cannot} express the
latter --- a near-zero excitation is ``no effect,'' never inhibition. Firm
inhibition therefore moves into the mark factor as a cooldown covariate, and
the exp-link block Hawkes (Section~\ref{sec:blockhawkes}) exists as the
single-layer alternative that can carry both signs at once, at the price of
abandoning the linear mean-field identity behind counts-only MBPP calibration
($\E[g(Y)]\neq g(\E Y)$; M5).

\paragraph{Excitation is a bounded add-on, not the headline.} Our prior, which
the scenario analysis sharpens, is that macro/sector covariates carry most of
the predictable variation in deal flow and self-excitation is a modest,
short-horizon correction. The architecture treats excitation accordingly:
nonnegative, stability-constrained (row-sum bound $\Rightarrow$ spectral radius
$<1$), reported always with its effective radius, and challenged by a standing
covariates-only null (\S\ref{sec:design_bench}). In the synthetic rehearsals
the covariate block indeed dominated held-out gains and the weekly excitation
channel was small or even harmful once a persistent latent factor existed
(\code{exp25}, \code{exp27}) --- a mechanism we expect to survive on real data
in the form of ``risk-on/risk-off'' common factors.

\paragraph{Weekly counts vs event times is a real fork.} The weekly GLM is
convenient and aggregation-coherent; the event-time model is the only one that
can see sub-week timing. Which to deploy is a scenario question
(\S\ref{sec:design_scenarios}, S2): the decision variable is the ratio of
bucket width to the contagion kernel's half-life.

% ---------------------------------------------------------------------
\subsection{Parametrizing the mark factor: conditional choice over risk sets}
\label{sec:design_mark}

The mark factor is a conditional-choice model: softmax over the live pool with
score $q_{i,t}=(w_0+u_s)^\top Z_{i,t}+\eta_s C_{i,t}$
\eqref{eq:cox_score}. The defenses, each tied to a data property:

\begin{itemize}[leftmargin=1.6em]
\item \emph{Conditional likelihood} because labels are positive-only (M8): no
invented non-events, no class-imbalance machinery, and the sector-time baseline
cancels exactly as in Cox partial likelihood.
\item \emph{The funded pool} because covariates only exist for previously
funded firms (LOCF reality); never-funded firms cannot be scored on firm
features and are routed to the newcomer channel
(\S\ref{sec:design_newcomer}).
\item \emph{The exclusion rule} (drop the most recently funded firm of the
sector) encodes the one suppression we are certain of; ties break
deterministically; risk-set construction precedes state updates (audited, unit
tested).
\item \emph{A single gap feature carries inhibition.} M6 derives the implied
hazard-vs-gap curves under an indicator cooldown and under exponential
inhibition, and what a binned refit recovers in each case --- the
functional-form check that transfers to real data.
\item \emph{Likelihood family is a convenience choice among equals.} On
identical pools and events, conditional logit, Cox partial likelihood, and a
discrete hazard were statistically indistinguishable and reached the
true-parameter ceiling in rehearsal (\code{exp25}); an intensity-based
(Hawkes-recency) ranker underperforms except where its dynamics are the truth.
We therefore claim robustness of the \emph{family}, and put the real leverage
where the rehearsals located it: pool correctness and feature quality, not the
link function.
\end{itemize}

% ---------------------------------------------------------------------
\subsection{The newcomer channel: every design implemented}
\label{sec:design_newcomer}

Cold start is structural: $\sim$27\% of choice events in our synthetic worlds
are won by first-timers or untracked firms, and the same order of magnitude is
plausible in production. Four designs are implemented (\code{exp28}; M7), in
increasing structural commitment:
N1, a context-covariate alternative-specific utility
$q_{0,s,t}=b_0+\gamma^\top g_{s,t}$ on market context (sector heat, historical
first-financing share, macro, pool size), making the newcomer share
time-varying and calibratable;
N2, a representative-entrant model using the estimable at-first-funding feature
distribution, with the exact Gaussian log-MGF utility
$w^\top\mu_s+\tfrac12 w^\top\Sigma_s w+c_s$ --- M7 both derives it and proves
the pool-size term is unidentified from a free constant unless the entrant pool
varies in time (an identifiability statement we can make \emph{before} real
data);
N3, a nested two-step (binary newcomer-vs-incumbent, then incumbents-only
choice) whose joint likelihood factorizes exactly and whose incumbent weights
must match N1's --- a consistency check more than a competitor;
N4, the structural reading: first financings are the \emph{immigrant} stream of
the sector branching process, repeats are the offspring stream; splitting the
sector counts into the two channels and fitting them separately aligns the
outside option with Hawkes semantics --- the newcomer alternative \emph{is} the
immigrant channel seen from the mark factor.

The synthetic worlds also let us quantify a conflation no real dataset can:
``newcomer'' mixes true entrants with untracked incumbents, and the private
deal map prices that conflation in joint likelihood. On real data the split is
unobservable; the scenario map (S4) says what to expect instead.

% ---------------------------------------------------------------------
\subsection{The scenario map: identifiability and expected performance}
\label{sec:design_scenarios}

This is the section we will reach for when real data arrives. Each scenario
states: what is identifiable, what silently is not, where we expect the model
to perform, the diagnostic that distinguishes the scenarios, and the reporting
discipline that keeps us honest. The synthetic worlds are cited as the
mechanism check for each claim.

\paragraph{S1 --- What drives deal flow: covariates, contagion, or a latent
factor.}
The three drivers are near-indistinguishable in weekly counts: a persistent
unobserved common factor loads onto the lagged-count channel and masquerades as
contagion (rehearsed: fitted effective radius $2.4\times$ truth under a mild
latent factor; and at pooled-sector level even day-resolution timestamps could
not separate latent quality from excitation). \emph{Expectations:} fitted
excitation is an \emph{upper bound}, its off-diagonal entries are not
structurally interpretable, and only orderings (which sector is most
self-exciting; which firms are most suppressed) survive --- orderings were
recoverable even under misspecification (rank correlation $0.94$ for
firm-level inhibition in the day-resolution rehearsal). \emph{Diagnostics:}
add a lagged aggregate-market covariate and watch the fitted radius drop; if it
collapses, the ``contagion'' was a common factor. \emph{Reporting:} predictive
claims freely; structural excitation claims only as bounds with the radius
attached.

\paragraph{S2 --- Observation resolution vs kernel timescale.}
The branching ratio $\kappa$ is identified at any resolution; the decay
$\theta$ is identified only when the bucket width is small relative to the
kernel (M3 gives the Fisher-information collapse; the rehearsal located the
cliff: buckets $\le$ one kernel half-life are safe, $\ge4$ half-lives bias the
decay irrecoverably). \emph{Expectations for real data:} with day-resolution
deal dates and a contagion half-life plausibly of order days-to-weeks, the
event-time track can estimate timing; monthly aggregations cannot, and any
$\theta$ fitted from them is an artifact. Volume matters as much as
resolution: at a few thousand events per sector the $\kappa$--$\theta$ ridge
dominates before bucketing even enters --- pool sectors or fix $\theta$ to a
bank of interpretable half-lives and profile. \emph{Diagnostic:} refit at two
bucket widths; a stable $\hat\theta$ is informative, a drifting one means you
are on the ridge. \emph{Reporting:} state the (resolution, volume) pair next
to any timing claim.

\paragraph{S3 --- Risk-set observability.}
The single largest sensitivity in the whole architecture. With truthful
entry/exit the mark factor recovers true economics and sits at its theoretical
ceiling; with pools inferred from observables the fitted weights can
\emph{invert} --- recency becomes a large positive predictor because it proxies
the unobserved exit information, and in rehearsal \emph{no} observable exit
model (even one predicting true exit at AUC $0.88$) closed the gap.
\emph{Expectations:} on real data, where graduation/death labels are partial,
we expect measured ranking quality to sit strictly between the observed-pool
and clean-pool bounds, and \emph{fitted weights to be biased in the direction
of survival proxies}. \emph{Diagnostics:} weight signs that contradict
economic priors (capital raised strongly positive, cooldown weak) are the
signature of pool contamination, not a market insight. \emph{Reporting:}
metrics under both the naive pool and the best exit-labeled pool; the spread is
the price of the observation gap and must be quoted, not averaged away.

\paragraph{S4 --- Universe coverage and the newcomer share.}
With partial coverage, the newcomer channel absorbs two indistinguishable
populations (true entrants, untracked incumbents). \emph{Expectations:} the
context-ASC (N1) keeps the newcomer \emph{share} calibrated over time even
though the composition is unknowable; the entrant-quality decomposition (N2) is
estimable only for the tracked-entrant part; a rising newcomer share with
constant context is the diagnostic that coverage, not the market, changed.
\emph{Reporting:} newcomer-share reliability by quarter alongside any
incumbent ranking metric; never compare incumbent top-$k$ across periods with
different coverage.

\paragraph{S5 --- Volume regime.}
Sparse worlds (tenths of events per sector-week) leave excitation numerically
unidentifiable and favor adaptive baselines; the rehearsals showed a
three-line EWMA matching structural models wherever a latent level factor
existed, and penalized (L1) excitation with a stability bound as the only way
the GLM stays sane. Dense worlds progressively identify excitation and then
timing. \emph{Expectations:} on real sector panels (tens of sectors, years of
weeks) we are in the sparse-to-middle regime: covariates + small bounded
excitation, timing claims only from the event-time track.
\emph{Diagnostic:} if a trivial EWMA matches the structural model out of
sample, believe it --- it is telling you the identifiable content is a slowly
moving level.

\paragraph{S6 --- Feature measurement noise.}
Errors-in-variables attenuates structure long before it harms prediction
(rehearsed at up to 20\% relative noise: ranking essentially unchanged,
weights attenuated $\times0.88$, recovery correlations visibly degraded).
\emph{Expectations:} PitchBook-style noise (undisclosed sizes, estimated
employees) will barely move top-$k$ but will shrink every fitted coefficient
toward zero --- coefficients are lower bounds in magnitude.
\emph{Reporting:} predictive and structural claims split, with attenuation
acknowledged on the latter.

\paragraph{S7 --- Stability regime of the market.}
Two explosion mechanisms exist (M5): additive supercriticality
($\rho\ge1$), and the log-link raw-count feedback that is explosive even at
$\rho(b)\ll1$ --- the mechanism we measured as metastability (a $0.35$-radius
world tipping into a supercritical episode). \emph{Expectations:} a genuinely
``hot'' market (AI 2021-style) will push fitted excitation toward the
stability boundary, and \emph{simulation}, not estimation, is where it breaks.
The guards are layered precisely for this scenario: stability-constrained
fitting, baseline winsorization, risk-set caps, first-$K$ caps, refusal to
simulate above effective radius $0.95$. \emph{Diagnostic:} row sums pinned at
the stability bound mean the likelihood wants more feedback than stability
allows --- treat as a latent-factor or regime-change signal, not as contagion.

% ---------------------------------------------------------------------
\subsection{Pre-registered expectations for real data}
\label{sec:design_prereg}

Written before any real-data fit, so that the first back-test confirms or
falsifies cleanly:

\begin{enumerate}[leftmargin=1.6em]
\item Macro/sector covariates will carry the large majority of the held-out
gain over a seasonal baseline; weekly excitation will add little and may hurt
until stability-penalized (S1, S5).
\item The fitted effective radius will exceed any plausible true contagion;
adding a market-level factor covariate will reduce it materially (S1).
\item $\theta$ (timing) will be estimable only on the day-resolution event
stream, and only for sectors with sufficient volume; weekly refits will sit on
the $\kappa$--$\theta$ ridge (S2).
\item Conditional logit, Cox, and hazard variants of the mark factor will be
indistinguishable out of sample; feature and pool changes will dominate model
changes (S3, \code{exp25}).
\item Fitted mark weights will lean toward survival proxies relative to
economic priors, to a degree that shrinks as exit labeling improves (S3).
\item The newcomer share will be calibratable within a couple of points by the
context ASC; its composition will remain unidentifiable (S4, M7).
\item A gradient-boosted ranker at feature parity will land within a few
points of the conditional logit; a large gap would falsify the linear-utility
restriction and would be the single most informative failure we could observe
(\S\ref{sec:design_bench}).
\end{enumerate}

% ---------------------------------------------------------------------
\subsection{Benchmarks as standing falsification tests}
\label{sec:design_bench}

Because we cannot yet validate on real data, every restriction gets a
challenger that travels with the codebase:
\textbf{B1} (\code{exp26}): XGBoost on identical choice sets, at feature parity
(isolates the linear-utility restriction) and with a strictly point-in-time
kitchen sink (bounds the value of missing features); decision rule
pre-registered --- parity means the parametric mark factor stands, a gap means
features/nonlinearity, and the gap size is the budget for enrichment.
\textbf{B2} (\code{exp27}): the covariates-only NHPP null with the
decomposition seasonal $\to$ covariates $\to$ +excitation $\to$ +event-time
timing; excitation must earn its held-out nats or be dropped from the
production configuration (in rehearsal it earned them only at event-time
resolution --- weekly excitation was marginal-to-harmful, while the event-time
block recovered genuine timing gain on the day-resolution world).
\textbf{The alternatives audit} (\code{exp25}): the model-family robustness
result cited throughout, itself protocol-audited adversarially before its
numbers were trusted.

% ---------------------------------------------------------------------
\subsection{Analytical MLE first; learned operators as accelerators}
\label{sec:design_pino}

Every estimator in the pipeline has an exact, concave or profile-concave
likelihood (M2, M4): nothing in \emph{calibration} needs learning. Learned
operators (PINO) enter only as forward-solve accelerators for scenario stress
and walk-forward simulation --- rehearsed at sub-percent-to-few-percent
solve error and millisecond latency --- and the amortized \emph{inverse}
direction was tested and rejected (direct regression plateaus far above
differentiable-inversion error). This ordering is itself a design decision:
transparency and identifiability first, acceleration second, and no learned
component anywhere a likelihood is available.

% ---------------------------------------------------------------------
\subsection{Limitations and the production checklist}
\label{sec:design_limits}

\begin{enumerate}[leftmargin=1.6em]
\item \textbf{Risk-set truthfulness} is assumed by the v2 architecture and is
the dominant sensitivity (S3). Production effort should go to exit/graduation
labeling before any model refinement; metrics reported under both pools.
\item \textbf{Point-in-time discipline}: end-of-sample snapshot columns are
forbidden as features; every feature path is audited and unit-tested.
\item \textbf{Overdispersion}: Poisson intervals are anti-conservative at the
measured dispersions; quasi-Poisson corrections accompany all NLL tables and a
negative-binomial sector model is the natural upgrade.
\item \textbf{Multi-deal weeks}: repeated softmax approximates Plackett--Luce;
the approximation is bounded and audited, and becomes material only if
multi-deal sector-weeks dominate.
\item \textbf{Synthetic evidence}: all numbers cited here validate
\emph{mechanisms} under known truths. The scenario map and the pre-registered
expectations of \S\ref{sec:design_prereg} are the bridge: the first real
back-test should be scored against them, item by item.
\end{enumerate}

\paragraph{Summary of the argument.} The factorization
\eqref{eq:design_factorization} makes the two-stage architecture lossless; the
data reality (LOCF covariates, positive-only labels, churning partial universe,
chosen resolution) fixes the parametrization of each factor; identifiability
analysis decides what is estimated, what is fixed, and what is only ever a
bound; the scenario map says in advance where the machinery will be strong
(prediction, orderings, calibrated shares) and where it will be silent or
biased (structural excitation under latent factors, timing under coarse
buckets, weights under contaminated pools); and every remaining discretionary
choice ships with a falsification test. What survives is deliberately
conservative: a covariate-driven sector intensity with small, bounded,
stability-constrained excitation, and a conditional risk-set choice model with
a context-driven outside option --- with the honest map of when to trust each
number attached.
